\section{A self-contained collision interface}
\label{sec:physical-bridge}
We give the preparation, mechanical estimate, observation law and
private simulator class needed for the physical specialization. The
following argument does not appeal to an unpublished physical lemma.
It retains the same positive-volume preparation and positive sensor
noise as the earlier collision construction.

\begin{definition}[The reset row--event architecture]
\label{def:physical-class}
A preparation supplies two sequential report slots and, afterwards, an
actual mark. The source reports in the two slots are constant. A private
width-one simulator has no retained information between those slots
and no persistent hidden seed. It may use fresh independent coins.
Its first emission is $X_1\sim\Ber(p)$; after receiving the current
gate $A\in\{0,1\}$, its second emission has a fixed Bernoulli kernel
with parameter $q_A$. The parameters $(p,q_0,q_1)$ are fixed throughout
the audit. The simulator does not observe the actual mark. All its
coins are independent of that mark. Different preparations reset all
unretained variables and are independent conditional on the parameters.

The auditor observes the first report, chooses a gate using its retained
state and that report, observes the second report, and then observes the
mark. During training it may choose rows and stop adaptively, using at
most the declared number of candidate preparations. It next freezes
one gate policy and one event, with any retained random choice included
in its register. It uses that same row--event continuation on one fresh
candidate preparation and one independent target preparation and returns
the difference of their event indicators, target minus candidate.
It may not reselect the event after either validation. Independent
target-only calibration is allowed but all retained calibration and
randomness are charged. Peak width counts every bit, mark, event and
validation-accumulator cut. Clocked schedules are declared externally;
an autonomous implementation must retain its phase. An unretained
transcript, seed or stopping-time value is not available to a continuation.
\end{definition}

In the all-on row the complete private law is $b(p,q_1)$, with an
independent fair actual mark. For lower bounds it suffices to restrict
the adversary to $q_0=0$, $q_1=q$. In this subfamily every deterministic
gate policy is the pushforward
\[
 T_g(x_1,x_2,w)=(x_1,g(x_1)x_2,w)
\]
of $b(p,q)$. The gate, when included in the transcript, is another
function of these same coordinates. Randomized gate policies are
covered by averaging their independent random tables. This proves
both the all-on product characterization for the upper construction
and the common-garbling subfamily used for the outer converse. The
width-one constraint is a simulator constraint; it is not an uncharged
restriction on the auditor's register.

\subsection{Preparation and elastic scattering}
Consider two labeled equal-mass hard spheres in $\mathbb R^3$, with no
walls and contact distance $d$. Write $r=q_1-q_2$, $g=v_1-v_2$,
$V=(v_1+v_2)/2$, and let the center of mass position range over any fixed
rational box of positive volume. Prepare $(B,W)\in\{0,1\}^2$ with
probability $5/16$ at $B=W$ and $3/16$ at $B\ne W$. Conditional on
$(B,W)$, use the uniform product preparation on the following boxes:
\begin{align}
 r&=(-a_0,(2B-1)b_0,c_0),&
 a_0&\in[29/10,31/10],& b_0&\in[9/20,11/20],\notag\\
 |c_0|&\le1/100,&g_x&\in[19/10,21/10],&
 |g_y|,|g_z|&\le1/100,\label{eq:local-packets}\\
 |V_x|,|V_y|&\le1/100,&&&|V_z-(2W-1)|&\le1/100.\notag
\end{align}
Thus the actual mark is $W=\mathbf1_{\{V_z>0\}}$, preserved by momentum
conservation. Set
\[
 h(z)=\operatorname{sgn}(z)\min\{1,(|z|-1/4)_+\}.
\]
This is a bounded $1$-Lipschitz detector, identically zero on the incoming
transverse-velocity range.

\begin{lemma}[Uniform collision signal]\label{lem:local-scattering}
For every preparation in \eqref{eq:local-packets} and every
$9/10\le d\le11/10$, exactly one collision occurs and its time lies in
$(4/5,3/2)$. Afterwards
\[
             (2B-1)v_{1,y}>1/2,\qquad (2B-1)h(v_{1,y})>1/4.
\]
For $0<d\le2/5$ no collision occurs on $[0,3]$ and the detector is zero
throughout that interval.
\end{lemma}
\begin{proof}
Before contact, $r(t)=r+tg$. At $t\le4/5$ its $x$ coordinate is at
most $-29/10+(4/5)(21/10)=-61/50$, so its norm exceeds $11/10$.
At $t=3/2$ its absolute coordinate bounds are $1/4$, $113/200$ and
$1/40$, whose squared sum is less than $(9/10)^2$. The first contact
therefore occurs in the claimed interval. At the free zero of the
$x$ coordinate, whose time is at most $31/19$, the transverse norm is
less than $9/10$. Hence the entry contact has $n_x<0$, where
$n=r(t_c)/d$. Its remaining coordinates give
\[
 |n_x|^2\ge1-\frac{(113/200)^2+(1/40)^2}{(9/10)^2}>(77/100)^2,
 \qquad (2B-1)n_y\ge87/220.
\]
It follows that $-g\cdot n>(19/10)(77/100)-2/100=1443/1000$.
The elastic rule gives
\[
 v_1^+=V+g/2-(g\cdot n)n,\qquad g^+=g-2(g\cdot n)n.
\]
Consequently
\[
 (2B-1)v_{1,y}^+>-3/200+(1443/1000)(87/220)>1/2.
\]
The definition of $h$ implies the signal bound. Since
$r(t_c)\cdot g^+>0$, squared separation increases strictly afterwards;
there is no second collision. Finally, for $d\le2/5$ and $0\le t\le3$,
$|(r+tg)_y|\ge9/20-3/100=21/50>d$. There is no contact, and the
incoming $|v_{1,y}|\le3/200$ makes $h$ zero.
\end{proof}

\subsection{Positive noise and the actual mark}
Let $\mathcal B_t$ be Brownian motion independent of the preparation and
set
\[
 Y_t=\int_0^t h(v_{1,y}(s))\,ds+\sigma\mathcal B_t,
                  \qquad 0<\sigma\le1/24.
\]
The two potential all-on reports are
\[
 Z_1=\mathbf1_{\{2(Y_2-Y_{3/2})>0\}},\qquad
 Z_2=\mathbf1_{\{2(Y_3-Y_{5/2})>0\}}.
\]
The actual second report is $AZ_2$. The gate changes acquisition, not
the mechanical path. The continuous path is implemented by the target
sensor; it is not made available to the private simulator or supplied
as an uncharged auditor buffer. Denote the law of $(Z_1,Z_2,W)$ by
$Q_{d,\sigma}$.

\begin{theorem}[Uniform physical presentation]\label{thm:local-physical}
For $9/10\le d\le11/10$ and $0<\sigma\le1/24$,
\begin{equation}\label{eq:local-tv}
                    \TV(Q_{d,\sigma},Q_*)<1/300.
\end{equation}
The same estimate holds after every declared feedback gate, including
its observable gate value. The all-on private product family and the
common-garbling subfamily in Definition~\ref{def:physical-class} apply
to this very experiment.
\end{theorem}
\begin{proof}
Every collision occurs before $3/2$. Conditional on the preparation,
the two normalized sensor windows are independent Gaussians with mean
$h(v_{1,y}^+)$ and standard deviation $\sqrt2\sigma\le1/12$.
Lemma~\ref{lem:local-scattering} makes their means have sign $2B-1$
and magnitude greater than $1/4$. Each potential report therefore
differs from $B$ with probability less than $\Phi(-3)$.
The Gaussian tail estimate
\[
 \Phi(-x)\le\frac{e^{-x^2/2}}{x\sqrt{2\pi}}
\]
gives $\Phi(-3)<1/600$: $e^{9/2}>80$ and $\sqrt{2\pi}>5/2$
suffice. Couple the ideal and noisy reports using the same preparation,
actual mark and gate random table. Except on a set of probability less
than $2\Phi(-3)<1/300$, both potential reports equal $B$ and the entire
gated transcript agrees. The ideal marked law is precisely $Q_*$,
by the prepared $(B,W)$ table. This proves \eqref{eq:local-tv} and the
feedback-uniform version. Finally, the source is blind, its actual mark
is unobserved by the simulator, and the gate is its only current input
at the second emission. The product and subfamily assertions follow
from the two independent private-emission kernels specified above.
\end{proof}

For $d\le2/5$, the two windows instead contain only independent centered
Gaussian noise. Their signs are independent fair bits, independent of
$W$, and a private width-one simulator reproduces every gated row
exactly. Thus the informative regime and its no-collision control use
the same detector, preparation convention and actual mark.
